\section{Podstawy przetwarzania danych}

\subsection{Metody redukcji wymiaru danych}

Redukcja wymiaru polega na zastąpieniu danych opisanych wieloma zmiennymi przez reprezentację o mniejszej liczbie wymiarów przy możliwie małej utracie istotnej informacji. Stosuje się ją dla wizualizacji, przyspieszenia algorytmów, usuwania szumu, walki z przekleństwem wymiarowości i poprawy generalizacji.

\subsubsection{Redukcja wymiaru a selekcja cech}

Redukcja wymiaru tworzy nowe zmienne będące transformacjami cech oryginalnych. Selekcja cech wybiera podzbiór istniejących cech. PCA tworzy nowe składowe liniowe; selekcja może wybrać np. tylko wiek, dochód i liczbę transakcji z oryginalnej tabeli.

\subsubsection{PCA}

Principal Component Analysis znajduje ortogonalne kierunki maksymalnej wariancji danych. Po centrowaniu macierzy danych $X$ szukamy kierunku $w$, który maksymalizuje:
\[
\operatorname{Var}(Xw)
\]
przy ograniczeniu $\|w\|=1$. Rozwiązaniem są wektory własne macierzy kowariancji:
\[
\Sigma=\frac{1}{n-1}X^TX.
\]
Pierwsza składowa główna wyjaśnia największą wariancję, druga największą wariancję w kierunku ortogonalnym do pierwszej itd.

Zalety PCA:
\begin{itemize}
  \item prostota,
  \item szybkość,
  \item dekorrelacja zmiennych,
  \item możliwość interpretacji przez wariancję wyjaśnioną.
\end{itemize}
Wady:
\begin{itemize}
  \item liniowość,
  \item wrażliwość na skalowanie cech,
  \item skupienie na wariancji, niekoniecznie na informacji predykcyjnej,
  \item ograniczona interpretowalność składowych.
\end{itemize}

\subsubsection{SVD}

Singular Value Decomposition rozkłada macierz:
\[
X=U\Sigma V^T.
\]
Przycięcie do $k$ największych wartości osobliwych daje najlepsze przybliżenie rzędu $k$ w normie Frobeniusa. SVD jest podstawą PCA i metod latent semantic analysis.

\subsubsection{LDA jako redukcja nadzorowana}

Linear Discriminant Analysis szuka kierunków dobrze rozdzielających klasy. Maksymalizuje stosunek wariancji międzyklasowej do wewnątrzklasowej. Dla $C$ klas redukuje do co najwyżej $C-1$ wymiarów. W przeciwieństwie do PCA używa etykiet.

\subsubsection{Metody nieliniowe}

\begin{description}
  \item[Kernel PCA] wykonuje PCA w przestrzeni cech zadanej przez kernel.
  \item[Isomap] zachowuje odległości geodezyjne na rozmaitości.
  \item[LLE] zakłada, że punkt można odtworzyć z lokalnych sąsiadów.
  \item[t-SNE] służy głównie do wizualizacji, zachowuje lokalne podobieństwa, ale zniekształca globalne odległości.
  \item[UMAP] szybka metoda wizualizacji i redukcji, często lepiej zachowuje strukturę globalną niż t-SNE, ale nadal wymaga ostrożnej interpretacji.
  \item[Autoenkodery] sieci neuronowe uczące kodu niskowymiarowego przez rekonstrukcję wejścia.
\end{description}

\subsubsection{Dobór liczby wymiarów}

Można użyć:
\begin{itemize}
  \item progu wariancji wyjaśnionej, np. 95\% w PCA,
  \item wykresu osypiska,
  \item walidacji krzyżowej dla jakości modelu predykcyjnego,
  \item metryk rekonstrukcji,
  \item stabilności reprezentacji.
\end{itemize}

\subsubsection{Co powiedzieć na obronie}

Redukcja wymiaru tworzy niskowymiarową reprezentację danych. PCA jest podstawową metodą liniową maksymalizującą wariancję, SVD daje jej stabilną implementację, LDA używa etykiet, a metody nieliniowe takie jak t-SNE, UMAP i autoenkodery wykrywają bardziej złożone struktury. Trzeba pamiętać o skalowaniu, utracie informacji i walidacji liczby wymiarów.

\subsection{Metody selekcji cech}

Selekcja cech wybiera podzbiór oryginalnych zmiennych. Celem jest poprawa generalizacji, redukcja kosztu pomiaru, przyspieszenie modelu, interpretowalność i zmniejszenie ryzyka overfittingu.

\subsubsection{Filtry}

Metody filtrujące oceniają cechy niezależnie od konkretnego modelu albo z użyciem prostych statystyk:
\begin{itemize}
  \item wariancja cechy,
  \item korelacja z etykietą,
  \item test $\chi^2$ dla cech kategorycznych,
  \item ANOVA F-test,
  \item mutual information,
  \item współczynnik korelacji rang Spearmana,
  \item usuwanie jednej z silnie skorelowanych cech.
\end{itemize}

Filtry są szybkie i skalowalne, ale mogą ignorować interakcje między cechami.

\subsubsection{Wrappery}

Wrapper używa modelu jako czarnej skrzynki i ocenia podzbiory cech przez walidację:
\begin{itemize}
  \item forward selection - zaczynamy od pustego zbioru i dodajemy cechy,
  \item backward elimination - zaczynamy od wszystkich i usuwamy,
  \item recursive feature elimination,
  \item przeszukiwanie heurystyczne, np. algorytm genetyczny.
\end{itemize}

Wrappery mogą dawać lepsze podzbiory dla konkretnego modelu, ale są kosztowne i podatne na przeuczenie, jeśli walidacja jest źle zorganizowana.

\subsubsection{Metody wbudowane}

Embedded methods wybierają cechy w trakcie uczenia modelu:
\begin{itemize}
  \item Lasso, czyli regularyzacja $L_1$, zeruje część współczynników,
  \item Elastic Net łączy $L_1$ i $L_2$,
  \item drzewa i lasy losowe dają importance cech,
  \item gradient boosting ma miary gain/split importance,
  \item modele liniowe z karami grupowymi wybierają grupy cech.
\end{itemize}

Trzeba uważać na interpretację importance w drzewach: cechy o wielu możliwych podziałach albo silnie skorelowane mogą być faworyzowane lub dzielić znaczenie.

\subsubsection{Stability selection}

Stability selection wielokrotnie losuje podpróbki danych, wykonuje selekcję i wybiera cechy pojawiające się stabilnie. Pomaga odróżnić cechy rzeczywiście istotne od przypadkowych, zwłaszcza gdy liczba cech jest duża.

\subsubsection{Selekcja a przeciek danych}

Selekcja cech musi być wykonywana wewnątrz pętli walidacji. Błąd: wybrać cechy na całym zbiorze, a potem zrobić cross-validation. Informacja ze zbioru testowego przecieka wtedy do modelu i wynik jest zawyżony.

\subsubsection{Co powiedzieć na obronie}

Selekcja cech wybiera podzbiór oryginalnych zmiennych. Filtry są szybkie i niezależne od modelu, wrappery oceniają podzbiory przez model i walidację, a metody wbudowane wybierają cechy podczas uczenia, np. Lasso albo drzewa. Najważniejsze jest unikanie przecieku danych i ocena stabilności wyboru.

\subsection{Metody postępowania z brakami w danych}

Braki danych mogą zniekształcić analizę i modelowanie. Wybór metody zależy od mechanizmu braków, udziału braków, typu zmiennej i modelu docelowego.

\subsubsection{Mechanizmy braków}

\begin{description}
  \item[MCAR] Missing Completely At Random - brak niezależny od danych obserwowanych i nieobserwowanych. Usunięcie wierszy nie wprowadza biasu, ale zmniejsza moc statystyczną.
  \item[MAR] Missing At Random - brak zależy od danych obserwowanych, ale nie od brakującej wartości po uwzględnieniu obserwowanych zmiennych. Można go modelować.
  \item[MNAR] Missing Not At Random - brak zależy od samej brakującej wartości albo ukrytej zmiennej. Najtrudniejszy przypadek.
\end{description}

Przykład MNAR: osoby z bardzo wysokim dochodem częściej nie podają dochodu.

\subsubsection{Usuwanie danych}

\begin{itemize}
  \item listwise deletion - usunięcie wierszy z jakimkolwiek brakiem,
  \item pairwise deletion - użycie dostępnych par obserwacji dla danej statystyki,
  \item usuwanie cech z bardzo dużym udziałem braków.
\end{itemize}

Usuwanie jest proste, ale może wprowadzać bias i marnować dane, jeśli braki nie są MCAR.

\subsubsection{Prosta imputacja}

\begin{itemize}
  \item średnia albo mediana dla zmiennych liczbowych,
  \item dominanta dla zmiennych kategorycznych,
  \item stała wartość specjalna, np. "unknown",
  \item imputacja grupowa, np. mediana w obrębie segmentu.
\end{itemize}

Mediana jest odporniejsza na odstające wartości niż średnia. Prosta imputacja zaniża wariancję i może osłabić zależności między zmiennymi.

\subsubsection{Imputacja modelowa}

\begin{itemize}
  \item k-nearest neighbors imputation,
  \item regresja predykcyjna,
  \item drzewa i lasy,
  \item EM dla modeli probabilistycznych,
  \item multiple imputation,
  \item autoenkodery albo modele generatywne.
\end{itemize}

Multiple imputation tworzy wiele uzupełnionych zbiorów, analizuje każdy i łączy wyniki, uwzględniając niepewność imputacji.

\subsubsection{Indykatory braków}

Czasem sam fakt braku jest informacyjny. Dodaje się cechę binarną:
\[
m_j =
\begin{cases}
1, & x_j\text{ było brakujące},\\
0, & \text{w przeciwnym razie}.
\end{cases}
\]
To bywa skuteczne przy MAR/MNAR, ale nie zastępuje analizy mechanizmu braków.

\subsubsection{Modele obsługujące braki}

Niektóre modele potrafią obsługiwać braki bez jawnej imputacji, np. wybrane implementacje drzew gradientowych mają kierunek domyślny dla braków. Nadal trzeba rozumieć, co model robi z brakami i czy zachowanie jest zgodne z problemem.

\subsubsection{Co powiedzieć na obronie}

Najpierw trzeba rozpoznać mechanizm braków: MCAR, MAR albo MNAR. Można usuwać wiersze lub cechy, imputować prostymi statystykami, imputować modelowo, stosować multiple imputation albo dodać indykatory braków. Imputację trzeba dopasować tylko na danych treningowych, aby uniknąć przecieku.

\subsection{Estymatory błędów}

Estymator błędu ocenia, jak dobrze model będzie działał na nowych danych. Błąd treningowy jest zwykle zbyt optymistyczny, bo model był dopasowany do tych danych.

\subsubsection{Podział train/validation/test}

Najprostszy schemat:
\begin{itemize}
  \item zbiór treningowy - uczenie parametrów,
  \item zbiór walidacyjny - dobór hiperparametrów i wybór modelu,
  \item zbiór testowy - końcowa bezstronna ocena.
\end{itemize}

Zbiór testowy powinien być użyty tylko raz na końcu. Wielokrotne podejmowanie decyzji na podstawie testu prowadzi do przeuczenia do testu.

\subsubsection{Walidacja krzyżowa}

W $k$-fold cross-validation dane dzieli się na $k$ części. Model trenuje się $k$ razy, za każdym razem na $k-1$ częściach, a testuje na pozostałej. Wyniki uśrednia się:
\[
\hat E_{CV}=\frac{1}{k}\sum_{i=1}^{k}E_i.
\]
Typowe wartości to $k=5$ albo $k=10$.

Leave-one-out CV to przypadek $k=n$. Ma mały bias, ale może mieć dużą wariancję i jest kosztowna.

\subsubsection{Stratyfikacja i dane zależne}

Dla klasyfikacji niezbalansowanej stosuje się stratified CV, aby proporcje klas były podobne w foldach. Dla szeregów czasowych nie wolno mieszać przyszłości z przeszłością; używa się walidacji kroczącej. Dla danych grupowych trzeba dzielić po grupach, aby obserwacje tej samej osoby, klienta albo urządzenia nie trafiały jednocześnie do treningu i testu.

\subsubsection{Bootstrap}

Bootstrap losuje próbki z powtórzeniami z danych treningowych. Obserwacje niewylosowane tworzą out-of-bag set. Bootstrap pozwala oceniać niepewność estymatora i stabilność modelu. W lasach losowych OOB error jest naturalnym estymatorem błędu.

\subsubsection{Macierz pomyłek i miary klasyfikacji}

Dla klasyfikacji binarnej:
\begin{itemize}
  \item TP - prawdziwie pozytywne,
  \item TN - prawdziwie negatywne,
  \item FP - fałszywie pozytywne,
  \item FN - fałszywie negatywne.
\end{itemize}

Miary:
\[
accuracy=\frac{TP+TN}{TP+TN+FP+FN},
\]
\[
precision=\frac{TP}{TP+FP},
\]
\[
recall=\frac{TP}{TP+FN},
\]
\[
F1=\frac{2\cdot precision\cdot recall}{precision+recall}.
\]
Accuracy jest mylące przy silnie niezbalansowanych klasach.

\subsubsection{Miary regresji}

\[
MAE=\frac{1}{n}\sum_i |y_i-\hat y_i|,
\]
\[
MSE=\frac{1}{n}\sum_i (y_i-\hat y_i)^2,
\]
\[
RMSE=\sqrt{MSE}.
\]
MAE jest odporniejsze na outliery, MSE/RMSE mocniej karzą duże błędy. $R^2$ mierzy część wariancji wyjaśnionej przez model.

\subsubsection{Bias i wariancja estymatora błędu}

Estymator błędu może być obciążony i mieć wariancję. Mały zbiór testowy daje dużą niepewność. Cross-validation zmniejsza zależność od pojedynczego podziału, ale wyniki foldów nie są całkowicie niezależne. Przy porównywaniu modeli warto stosować testy statystyczne albo przedziały ufności.

\subsubsection{Co powiedzieć na obronie}

Błąd modelu trzeba estymować na danych nieużytych do uczenia. Podstawowe metody to holdout, train/validation/test, walidacja krzyżowa, bootstrap i OOB. Miara błędu zależy od zadania: accuracy, precision, recall, F1, AUC dla klasyfikacji; MAE, MSE, RMSE dla regresji. Należy unikać przecieku danych i dobierać schemat walidacji do struktury danych.

\section{Uczenie ze wzmocnieniem}

\subsection{Uczenie ze wzmocnieniem: pasywne, aktywne i polityki}

Uczenie ze wzmocnieniem bada sytuację, w której agent uczy się przez interakcję ze środowiskiem. Agent wykonuje akcje, obserwuje następne stany i otrzymuje nagrody. Nie dostaje poprawnych etykiet akcji, lecz sygnał oceny skutków działania.

\subsubsection{Cykl interakcji}

W chwili $t$ agent jest w stanie $S_t$, wybiera akcję $A_t$, środowisko zwraca nagrodę $R_{t+1}$ i nowy stan $S_{t+1}$. Celem jest maksymalizacja oczekiwanego zwrotu:
\[
G_t=R_{t+1}+\gamma R_{t+2}+\gamma^2R_{t+3}+\dots
=\sum_{k=0}^{\infty}\gamma^k R_{t+k+1},
\]
gdzie $\gamma\in[0,1]$ jest współczynnikiem dyskontowym. Dla zadań epizodycznych suma kończy się w stanie terminalnym.

\subsubsection{Polityka}

Polityka opisuje sposób wyboru akcji:
\[
\pi(a|s)=P(A_t=a\mid S_t=s).
\]
Polityka deterministyczna przypisuje każdemu stanowi jedną akcję:
\[
\pi(s)=a.
\]
Polityka stochastyczna przypisuje rozkład prawdopodobieństwa na akcjach. Stochastyczność jest ważna dla eksploracji i dla środowisk częściowo obserwowalnych albo wieloagentowych.

\subsubsection{Pasywne uczenie ze wzmocnieniem}

W pasywnym RL polityka jest dana. Agent nie decyduje, jak ją ulepszyć, lecz ocenia jej jakość. Typowe zadanie: policy evaluation, czyli estymacja funkcji wartości $V^\pi(s)$ dla ustalonej polityki $\pi$.

Przykład: mamy strategię gry w blackjacka "dobieraj do 20, potem stop" i chcemy oszacować wartość stanów przy tej strategii. Metody: Monte Carlo prediction, TD(0), dynamic programming, jeśli model jest znany.

\subsubsection{Aktywne uczenie ze wzmocnieniem}

W aktywnym RL agent ma znaleźć dobrą albo optymalną politykę. Musi jednocześnie:
\begin{itemize}
  \item eksplorować, aby poznać skutki akcji,
  \item eksploatować, aby wybierać akcje już znane jako dobre.
\end{itemize}

To problem exploration-exploitation. Strategie:
\begin{itemize}
  \item $\eps$-greedy - z prawdopodobieństwem $1-\eps$ wybierz najlepszą znaną akcję, z $\eps$ losową,
  \item softmax/Boltzmann exploration,
  \item upper confidence bound,
  \item Thompson sampling,
  \item dodawanie szumu do parametrów albo akcji.
\end{itemize}

\subsubsection{On-policy i off-policy}

W uczeniu on-policy uczymy wartość albo poprawiamy tę samą politykę, która generuje dane. Przykład: SARSA.

W uczeniu off-policy dane generuje polityka behawioralna $b$, a uczymy politykę docelową $\pi$. Przykład: Q-learning. Off-policy pozwala uczyć politykę zachłanną na danych z polityki eksploracyjnej, ale może być mniej stabilne, zwłaszcza z aproksymacją funkcji.

\subsubsection{Funkcje wartości}

Wartość stanu:
\[
V^\pi(s)=\mathbb{E}_\pi[G_t\mid S_t=s].
\]
Wartość akcji:
\[
Q^\pi(s,a)=\mathbb{E}_\pi[G_t\mid S_t=s,A_t=a].
\]
Funkcja $V$ mówi, jak dobry jest stan przy danej polityce. Funkcja $Q$ mówi, jak dobra jest akcja w stanie, dlatego jest szczególnie użyteczna, gdy nie znamy modelu przejść.

\subsubsection{Co powiedzieć na obronie}

RL polega na uczeniu przez interakcję: stan, akcja, nagroda, kolejny stan. Polityka mapuje stany na akcje albo rozkłady akcji. Pasywne RL ocenia zadaną politykę, aktywne RL szuka polityki optymalnej i musi rozwiązać dylemat eksploracji i eksploatacji. Wartości $V^\pi$ i $Q^\pi$ opisują oczekiwany zwrot.

\subsection{Decyzyjny proces Markowa: definicja i przykłady}

MDP jest formalnym modelem sekwencyjnego podejmowania decyzji przy niepewności. Zakłada własność Markowa: przyszłość zależy od bieżącego stanu i akcji, a nie od pełnej historii.

\subsubsection{Definicja}

Skończony MDP można zapisać jako
\[
M=(S,A,p,r,\gamma),
\]
gdzie:
\begin{itemize}
  \item $S$ - zbiór stanów,
  \item $A$ - zbiór akcji,
  \item $p(s',r|s,a)$ - prawdopodobieństwo przejścia do stanu $s'$ i otrzymania nagrody $r$ po wykonaniu akcji $a$ w stanie $s$,
  \item $r$ - funkcja nagrody, często jako wartość oczekiwana,
  \item $\gamma\in[0,1]$ - współczynnik dyskontowy.
\end{itemize}

Warunek normalizacji:
\[
\sum_{s'\in S}\sum_{r\in R}p(s',r|s,a)=1
\quad\text{dla każdego }s,a.
\]

Z rozkładu $p(s',r|s,a)$ można wyznaczyć prawdopodobieństwo przejścia:
\[
p(s'|s,a)=\sum_r p(s',r|s,a),
\]
oraz oczekiwaną nagrodę:
\[
r(s,a)=\sum_r\sum_{s'} r\,p(s',r|s,a).
\]

\subsubsection{Własność Markowa}

Proces spełnia własność Markowa, jeśli:
\[
P(S_{t+1},R_{t+1}\mid S_0,A_0,\dots,S_t,A_t)
=P(S_{t+1},R_{t+1}\mid S_t,A_t).
\]
Stan powinien więc zawierać całą informację z historii potrzebną do przewidywania przyszłości. Jeśli obserwacja nie jest pełnym stanem, mamy POMDP, czyli częściowo obserwowalny MDP.

\subsubsection{Równania Bellmana}

Dla ustalonej polityki:
\[
V^\pi(s)=\sum_a\pi(a|s)\sum_{s',r}p(s',r|s,a)\left[r+\gamma V^\pi(s')\right].
\]
Dla funkcji akcji:
\[
Q^\pi(s,a)=\sum_{s',r}p(s',r|s,a)\left[r+\gamma\sum_{a'}\pi(a'|s')Q^\pi(s',a')\right].
\]

Optymalna funkcja wartości spełnia:
\[
V^*(s)=\max_a\sum_{s',r}p(s',r|s,a)\left[r+\gamma V^*(s')\right].
\]
Optymalna funkcja akcji:
\[
Q^*(s,a)=\sum_{s',r}p(s',r|s,a)\left[r+\gamma\max_{a'}Q^*(s',a')\right].
\]

\subsubsection{Przykłady MDP}

\begin{description}
  \item[Recycling robot] stan to poziom baterii, np. high/low. Akcje to search, wait, recharge. Nagrody zależą od zebrania śmieci i kar za rozładowanie.
  \item[Gridworld] stan to pozycja agenta na kratownicy, akcje to ruchy, przejścia mogą być stochastyczne, a nagrody przypisane do pól terminalnych i kosztów ruchu.
  \item[Blackjack] stan opisuje sumę kart gracza, odkrytą kartę krupiera i posiadanie używalnego asa. Akcje to hit/stand, nagroda końcowa to wygrana, przegrana albo remis.
  \item[Cart-pole] stan obejmuje pozycję wózka, prędkość, kąt słupka i prędkość kątową. Akcje to siła w lewo/prawo, nagroda za utrzymanie równowagi.
  \item[Bioreaktor] stan to odczyty sensorów i skład, akcje to ustawienia temperatury i mieszania, nagroda to tempo produkcji pożądanej substancji.
\end{description}

\subsubsection{Planowanie i uczenie}

Jeśli model $p$ i $r$ jest znany, możemy planować metodami programowania dynamicznego: policy evaluation, policy iteration, value iteration. Jeśli model nie jest znany, agent uczy się z doświadczenia metodami Monte Carlo, TD, SARSA, Q-learning albo metodami z aproksymacją funkcji.

\subsubsection{Co powiedzieć na obronie}

MDP to model $(S,A,p,r,\gamma)$, w którym przyszłość zależy tylko od bieżącego stanu i akcji. Polityka wybiera akcje, funkcje $V^\pi$ i $Q^\pi$ opisują oczekiwany zwrot, a równania Bellmana wiążą wartość stanu z nagrodą i wartością następnych stanów. Przykłady to robot sprzątający, gridworld, blackjack, cart-pole i sterowanie procesem.

\subsection{Metoda różnic czasowych TD-learning: idea, zakres stosowania, wady i zalety oraz porównanie z Monte Carlo}

Temporal Difference learning łączy idee Monte Carlo i programowania dynamicznego. Tak jak Monte Carlo uczy się z doświadczenia, bez pełnego modelu środowiska. Tak jak programowanie dynamiczne używa bootstrappingu, czyli aktualizuje wartość na podstawie obecnego oszacowania wartości następnego stanu.

\subsubsection{TD(0)}

Dla ustalonej polityki $\pi$ aktualizacja TD(0) ma postać:
\[
V(S_t)\leftarrow V(S_t)+\alpha\left[R_{t+1}+\gamma V(S_{t+1})-V(S_t)\right].
\]
Wyrażenie
\[
\delta_t=R_{t+1}+\gamma V(S_{t+1})-V(S_t)
\]
to błąd TD.

Cel aktualizacji TD to:
\[
R_{t+1}+\gamma V(S_{t+1}),
\]
czyli jednkrokowe przybliżenie zwrotu. W Monte Carlo celem jest pełny zwrot $G_t$ po zakończeniu epizodu.

\subsubsection{Algorytm predykcji TD(0)}

\begin{enumerate}
  \item Zainicjalizuj $V(s)$ dowolnie.
  \item Dla każdego epizodu:
  \begin{enumerate}
    \item rozpocznij w stanie $S$,
    \item wybierz akcję zgodnie z polityką $\pi$,
    \item wykonaj akcję, obserwuj $R$ i $S'$,
    \item zaktualizuj
    \[
    V(S)\leftarrow V(S)+\alpha\left[R+\gamma V(S')-V(S)\right],
    \]
    \item ustaw $S\leftarrow S'$,
    \item powtarzaj do stanu terminalnego.
  \end{enumerate}
\end{enumerate}

\subsubsection{Porównanie TD i Monte Carlo}

\begin{longtable}{p{0.24\textwidth}p{0.32\textwidth}p{0.32\textwidth}}
\toprule
Cecha & Monte Carlo & TD \\
\midrule
Model środowiska & nie wymaga & nie wymaga \\
Moment aktualizacji & po epizodzie & po każdym kroku \\
Cel aktualizacji & pełny zwrot $G_t$ & $R_{t+1}+\gamma V(S_{t+1})$ \\
Epizodyczność & wymaga epizodów albo specjalnych wariantów & działa w zadaniach ciągłych \\
Bootstrapping & nie & tak \\
Bias & mniejszy przy poprawnych próbkach & większy przez bootstrapping \\
Wariancja & często większa & zwykle mniejsza \\
Szybkość uczenia online & wolniejsza informacja & szybka aktualizacja \\
\bottomrule
\end{longtable}

MC jest nieobciążone względem rzeczywistych zwrotów przy próbkowaniu z polityki, ale może mieć dużą wariancję. TD wprowadza bias, bo używa własnych oszacowań, ale zwykle ma mniejszą wariancję i uczy się szybciej w zadaniach sekwencyjnych.

\subsubsection{SARSA}

SARSA jest on-policy metodą kontroli TD:
\[
Q(S_t,A_t)\leftarrow Q(S_t,A_t)+\alpha\left[R_{t+1}+\gamma Q(S_{t+1},A_{t+1})-Q(S_t,A_t)\right].
\]
Nazwa pochodzi od sekwencji:
\[
S_t,A_t,R_{t+1},S_{t+1},A_{t+1}.
\]
Ponieważ używa akcji faktycznie wybranej przez bieżącą politykę, uczy wartości polityki eksploracyjnej.

\subsubsection{Q-learning}

Q-learning jest off-policy:
\[
Q(S_t,A_t)\leftarrow Q(S_t,A_t)+\alpha\left[R_{t+1}+\gamma\max_a Q(S_{t+1},a)-Q(S_t,A_t)\right].
\]
Uczy wartości polityki zachłannej, nawet jeśli dane są generowane przez politykę eksploracyjną, np. $\eps$-greedy.

\subsubsection{Zakres stosowania}

TD jest użyteczne, gdy:
\begin{itemize}
  \item środowisko jest nieznane,
  \item dane przychodzą online,
  \item epizody są długie albo nie kończą się,
  \item chcemy aktualizować po każdym kroku,
  \item pełny zwrot Monte Carlo miałby dużą wariancję.
\end{itemize}

TD jest podstawą wielu metod RL, także z aproksymacją funkcji i deep RL.

\subsubsection{Wady}

\begin{itemize}
  \item bias przez bootstrapping,
  \item wrażliwość na krok uczenia $\alpha$,
  \item możliwa niestabilność przy aproksymacji funkcji, off-policy i bootstrappingu,
  \item potrzeba eksploracji dla metod kontroli,
  \item trudność doboru reprezentacji stanu.
\end{itemize}

\subsubsection{Zalety}

\begin{itemize}
  \item uczenie online,
  \item brak potrzeby modelu środowiska,
  \item brak potrzeby czekania do końca epizodu,
  \item często mniejsza wariancja niż MC,
  \item możliwość stosowania w zadaniach ciągłych,
  \item naturalne połączenie z funkcjami wartości akcji.
\end{itemize}

\subsubsection{Co powiedzieć na obronie}

TD-learning aktualizuje wartość stanu po każdym kroku, używając różnicy między starą wartością a celem $R+\gamma V(S')$. Błąd TD to $\delta=R+\gamma V(S')-V(S)$. W porównaniu z Monte Carlo TD nie czeka do końca epizodu i zwykle ma mniejszą wariancję, ale jest obciążone przez bootstrapping. SARSA jest on-policy, Q-learning off-policy.
